\section*{Abstract}

The 4D torus is not a postulate of T0 theory. It follows necessarily from three independent conditions: (1) freedom from divergences in quantum field theory requires periodic topology, (2) the value $\xipar = 4/30000$ encodes exactly four spacetime dimensions, (3) stable winding numbers on the torus generate the observed quantum numbers. The minimal length scale $L_0 = \xipar \cdot \lP$ is geometrically the length of a single torus lattice cell — the smallest stable winding. The Lagrangian and the torus yield $L_0$ consistently via two independent paths.

\section{Why a Fourth Dimension?}

\subsection{The Divergence Problem of Quantum Field Theory}

Quantum field theory suffers from a fundamental problem: the vacuum energy diverges:
\begin{equation}
	E_{\text{vac}} = \sum_n \frac{1}{2}\hbar\omega_n \rightarrow \infty
\end{equation}
The observed cosmological constant is $10^{120}$ times smaller than the QFT prediction. The Standard Model resolves this through artificial renormalisation. T0/FFGFT resolves it geometrically.

\subsection{Periodic Geometry as Solution}

On an infinite line the harmonic series diverges. On a circle it converges — Euler 1735:
\begin{equation}
	\sum_{n=1}^{\infty} \frac{1}{n^2} = \frac{\pi^2}{6}
\end{equation}

\begin{keypoint}[Topological Conclusion]
The periodic geometry automatically suppresses high modes. If spacetime has a periodic topology, UV divergences vanish geometrically — without artificial renormalisation. The $\pi^2$ in Euler's result is not number theory, but the signature of the topology.
\end{keypoint}

\subsection{Four Dimensions from \texorpdfstring{$\xi$}{xi}}

The parameter $\xipar = 4/30000$ encodes the dimensionality of spacetime:
\begin{equation}
	\xipar_d \sim (10^{-1})^d \quad \Rightarrow \quad \text{for } d=4: \quad \xipar_4 \sim 10^{-4}
\end{equation}

\begin{keyresult}[Four Dimensions from \texorpdfstring{$\xi$}{xi}]
The \textbf{4} in the numerator of $\xipar = 4/30000$ stands directly for the four spacetime dimensions. The fourth dimension is not postulated — it is already contained in the value of $\xipar$.
\end{keyresult}

\section{Why a Torus?}

\subsection{Winding Numbers as Conserved Quantities}

On a torus, winding numbers are topologically conserved. They generate stable discrete states:
\begin{align}
	\text{Spin} &= \text{winding number } w = n_\phi / n_\theta \\
	\text{Charge} &= \text{flux } \Phi = n \cdot h/e \\
	\text{Masses} &= \text{torus overtone frequencies}
\end{align}

\begin{keypoint}[Torus vs.\ Sphere]
A sphere $S^3$ is simply connected — no stable winding numbers. The torus is not simply connected — it generates exactly the observed quantum numbers. This is the topological reason for choosing the torus.
\end{keypoint}

\subsection{Torus Eigenvalues and Particle Masses}

The Laplace operator on a $d$-dimensional torus with radii $R_i$ has eigenvalues:
\begin{equation}
	E_{\vec{n}} = \frac{\hbar^2}{2} \sum_i \frac{n_i^2}{R_i^2}, \qquad \vec{n} \in \mathbb{Z}^d
\end{equation}
Particle masses appear as torus overtone frequencies with step size $1/3 = 1/D_f$:
\begin{equation}
	m_n = m_0 \cdot n^{1/3}
\end{equation}
The Koide formula for lepton masses follows as a geometric consequence.

\subsection{Compactification}

The fourth dimension is rolled up to radius $r_4$. This radius must coincide with the sub-Planck scale from the Lagrangian — otherwise the theory would be inconsistent. The condition:
\begin{equation}
	r_4 = L_0 = \xipar \cdot \lP
\end{equation}
The topology is therefore:
\begin{equation}
	\mathbb{R}^3 \times S^1 \quad \text{with} \quad r_4 = \xipar \cdot \lP
\end{equation}

\section{Derivation of \texorpdfstring{$L_0$}{L0} from Torus Geometry}

\subsection{The Winding Density}

The torus has $f = 1/\xipar$ sub-Planck cells per Planck length. With $\xipar = 4/30000$:
\begin{equation}
	f = \frac{1}{\xipar} = \frac{30000}{4} = 7500 \quad \text{cells per } \lP
\end{equation}
This number is not arbitrary — it follows directly from the geometric winding density $30000/4$ of the torus. The prime factorisation $7500 = 2^2 \times 3 \times 5^4$ encodes the symmetry structure of spacetime.

\subsection{Smallest Stable Winding}

The smallest stable winding on the torus has the length of a single lattice cell:
\begin{equation}
	L_0 = \frac{\lP}{f} = \lP \cdot \xipar
\end{equation}
Numerically:
\begin{equation}
	L_0 = 1{.}616 \times 10^{-35}\,\text{m} \times \frac{4}{30000} = 1{.}616 \times 10^{-35}\,\text{m} \times 1{.}333 \times 10^{-4}
\end{equation}

\begin{keyresult}[Minimal Length Scale from Torus Geometry]
\begin{equation}
	\boxed{L_0 = \xipar \cdot \lP = 2{.}155 \times 10^{-39}\,\text{m}}
\end{equation}
$L_0$ is the length of a single torus lattice cell — the smallest stable winding.
\end{keyresult}

\subsection{Geometric Interpretation}

$L_0$ is not simply $\xipar$ times the Planck length — $L_0$ is the length of a single torus lattice cell. Below $L_0$:
\begin{itemize}
	\item The concept of ``distance'' loses its physical meaning.
	\item All vacuum fluctuations are fully effective.
	\item The fractal correction $\Kfrak$ becomes structurally necessary.
\end{itemize}

\subsection{Consistency of the Two Derivation Paths}

$L_0$ is not derived twice — the Lagrangian and the torus yield the same number via two independent paths:

\begin{keyresult}[Two Consistent Paths to \texorpdfstring{$L_0$}{L0}]
\begin{align}
	\text{Lagrangian:} \quad & g_T = \xipar \cdot m \;\Rightarrow\; \xipar = 4/30000 \notag \\
	& \Downarrow \notag \\
	\text{Torus:} \quad & f = 1/\xipar = 7500 \text{ cells per } \lP \notag \\
	& \Downarrow \notag \\
	\text{Smallest winding:} \quad & L_0 = \lP / f = \xipar \cdot \lP \notag \\
	& \Downarrow \notag \\
	& \boxed{L_0 = 2{.}155 \times 10^{-39}\,\text{m}} \quad \text{(consistent from both paths)}
\end{align}
The Lagrangian provides $\xipar$. The torus explains what $\xipar$ means geometrically: the reciprocal winding density.
\end{keyresult}

\section{Comparison with Other Approaches}

\subsection{String Theory}

String theory requires 10 or 11 dimensions. The T0 Lagrangian with a single parameter $\xipar$ suffices with 4. Every additional dimension would be a postulate — in FFGFT there are no free parameters.

\subsection{Loop Quantum Gravity}

LQG discretises space into spin networks. T0 discretises spacetime through the periodic torus geometry. The difference: in T0 the discreteness is a topological consequence, not an additional postulate.

\subsection{Asymptotic Safety}

AS seeks a UV fixed point as the result of renormalisation group flow. In T0 there is no renormalisation group flow — $\xipar$ is a geometric constant, not a running parameter. The fixed point is not sought but established geometrically from the outset.

\section{Physical Consequences}

\subsection{Vacuum Catastrophe Resolved}

On the fractal torus, the vacuum energy converges through periodic suppression of high frequencies. No artificial renormalisation required.

\subsection{Three Generations as Overtone Structure}

Electron, muon and tau are not three independent particles — they are the fundamental tone and first overtones of a single torus resonator. The hierarchy follows from the torus eigenvalues:
\begin{equation}
	m_e : m_\mu : m_\tau \approx 1 : 206{.}8 : 3477
\end{equation}

\subsection{No Big Bang and No Frequency Dependence}

The torus has no boundaries. The observed redshift does not arise through expansion but through geometric energy loss — fractal summation extends the effective path length of photons through the T0 vacuum field:
\begin{equation}
	z \approx \xipar \cdot \ln\!\left(\frac{d}{\lP}\right)
\end{equation}

\begin{keypoint}[Frequency Independence of the Redshift]
Since the path extension is a property of the spacetime geometry itself, red and blue photons starting from the same location take exactly the same extended path. Their wavelengths are therefore stretched by the same percentage — there is no frequency dependence. This distinguishes the T0 mechanism fundamentally from simple ``tired light'' models, which fail precisely on the observed frequency independence. Derived in detail in Document 026 (Geometric Cosmology).
\end{keypoint}

\section{Summary}

\begin{keyresult}[Three Independent Conditions Point to the Same Torus]
\begin{align}
	(1)\; & \text{Freedom from divergences} \;\Rightarrow\; \text{periodic topology} \notag \\
	(2)\; & \xipar = 4/30000 \;\Rightarrow\; \text{exactly 4 dimensions} \notag \\
	(3)\; & \text{Winding numbers} \;\Rightarrow\; \text{stable quantum numbers} \notag
\end{align}
All three point to the same conclusion:
\begin{equation}
	\boxed{\mathbb{R}^3 \times S^1 \quad \text{with} \quad r_4 = L_0 = \xipar \cdot \lP = 2{.}155 \times 10^{-39}\,\text{m}}
\end{equation}
\end{keyresult}

\begin{foundation}[Note: Connection to the SI System]
The absolute numerical values in this document — in particular $L_0 = 2{.}155 \times 10^{-39}\,\text{m}$ — presuppose a connection to the SI unit system. This is not circular but epistemologically necessary: from a dimensionless parameter $\xipar$ alone, no absolute metre values can be derived. The connection is made via $m_e$, $\lP$ or $\alpha$ as equivalent entry points. Treated in detail in Document 180 (Section: Connection to the SI System) and in Documents 056 and 101.

Documents 180, 181 and 182 form a unit and should not be read separately.
\end{foundation}
